% =====================================================================
%  Chapter 2 — The Two Axes of the Hazard Scale
% =====================================================================

\chapter{The Two Axes of the Hazard Scale}
\label{ch:two-axes}

The Torino Impact Hazard Scale, adopted by the International Astronomical Union
in 1999, is the instrument the planetary-defence community uses to say how
much attention an object deserves. It is worth reading carefully, because the
reading is where the accountability gap becomes visible.

The scale is defined on two axes. The horizontal axis is the probability
\(p\) that the object will strike Earth. The vertical axis is the kinetic
energy \(E\) the object would release, measured in megatons of TNT. Every
assignment of a level to an object is a statement about the pair \((p, E)\),
and the pair matters more than either number alone: an object with a
probability of one in ten thousand and an energy of a kiloton is a different
problem from an object with the same probability and an energy of a hundred
megatons.

\section{Why the energy axis is the hard one}

The probability axis is measured. An orbit is fitted from astrometry, the
orbit is propagated, and the probability of intersection with the Earth's
cross-section is computed. The machinery is good, the observation times are
dense, and the numerical work is unremarkable. The energy axis is inferred,
and the inference has two stages.

The first stage is the diameter. A celestial body's absolute magnitude \(H\)
and its geometric albedo \(p\) are related to its effective diameter \(D\) by
the standard relation
\begin{equation}
D = \frac{1329}{\sqrt{p}}\,10^{-H/5}\ \ \mathrm{km},
\label{eq:hdp}
\end{equation}
which is a statement that a brighter object of the same albedo is bigger.
The problem is that \(p\) is a composition measurement, and the small
near-Earth objects that matter most have had their composition measured for
about fifteen per cent of them.

The second stage is the density. The mass is \(m = \rho \pi D^{3}/6\) for a
sphere of the given diameter, and the kinetic energy at impact speed \(v\) is
\begin{equation}
E = \frac{1}{2}\,m\,v^{2} = \frac{1}{2}\,\rho\,\frac{\pi}{6}\,D^{3}\,v^{2}.
\label{eq:energy}
\end{equation}
For a near-Earth object at a typical impact speed of 14 to 17 kilometres per
second, the density\idx{Density} is the dominant uncertainty. A rubble pile with a bulk
density\idx{Density} of 1,300~kg\,m\(^{-3}\) and a coherent monolith at 3,500~kg\,m\(^{-3}\)
differ by a factor of 2.7 in mass at the same diameter, and a factor of 2.7 in
energy.

\section{The error budget, in one line}

Differentiating \eqref{eq:energy} in logarithms gives the relation the whole
book turns on:
\begin{equation}
\delta\ln E = 3\,\delta\ln D + \delta\ln \rho.
\label{eq:dlne}
\end{equation}
The diameter enters three times because volume goes as the cube. This has a
plain consequence. If the energy is to be known to within a factor of three
and the density\idx{Density} to within thirty per cent, then the diameter must be known to
within thirty per cent. If the energy is to be known to within a factor of
two, the diameter must be known to about fourteen per cent.

\begin{table}[!htbp]
\centering
\small
\caption{What diameter accuracy the energy chain demands. The density is held
at thirty per cent, which is the state of the art for the near-Earth class.}
\label{tab:diameter-need}
\begin{tabularx}{\textwidth}{@{}P{1.19in}P{1.37in}Y@{}}
\toprule
\rowcolor{panel}
\textbf{Energy known to} & \textbf{Diameter must be known to} &
\textbf{Where that stands} \\
\midrule
a factor of 3 & 30 per cent & the present state of the art \\
\rowcolor{panel}
a factor of 2 & 14 per cent & beyond the published state of the art \\
a factor of 1.5 & 7 per cent & well beyond it \\
\bottomrule
\end{tabularx}
\end{table}

The NASA Center for Near Earth Object Studies' impact monitoring service
states the position in its own terms: the tabulated impact energy\idx{Impact energy} is
``generally good to within a factor of three'', and it attributes this to mass
uncertainty dominating the budget. That is the honest characterisation of the
field's present position, and it is the number the rest of this book is trying
to move.

\section{The threshold that moves when the size moves}

The \idxw{Torino scale}'s Level 1 boundary is not a fixed probability. It is a line
in the \((p, E)\) plane, and the line's intercept depends on the energy:
\begin{equation}
p \geq \frac{10^{-7/3}}{E} = \frac{4.64\times 10^{-3}}{E}\qquad (E \text{ in Mt}).
\label{eq:level1}
\end{equation}
This follows directly from the scale's definition: Level 1 is entered when
the combination \((\log_{10} E + \tfrac{1}{3} + \log_{10} p + 2)/2\) reaches
zero.

The consequence is the single most uncomfortable fact in this chapter. An
object whose astrometry has not changed at all --- not by one milliarcsecond,
not by one orbit --- can be Level 0 or Level 1 purely because of how it was
sized. A factor of 4.5 in energy moves the Level 1 probability threshold by
the same factor of 4.5, from about \(4.7\times 10^{-4}\) to about
\(2.2\times 10^{-3}\). Whether an object is a Level 1 event, and therefore
whether it generates a formal notification through the international warning
network, is a decision made by an assumption about the object's composition.

\section{What the scale does not do}

It is worth being clear about what the scale is for, because the gap this
book describes is not a failure of the scale. The Torino scale\idx{Torino scale} is a
communication instrument. It exists to let a duty officer, a news desk and a
commissioning agency talk about the same object in the same terms. It is
doing that job.

What it cannot do is supply the energy. The scale is handed a number and
reports a level; the number arrives from the same inference described above,
and the inference is where the uncertainty lives. The scale's own structure
makes the point: it is a two-axis system in which one axis is measured to
high precision and the other is assumed. Chapter 5 works through a case in
which the assumed axis moved by a factor of 4.5 while the measured axis stood
still.

\begin{keytakeaways}
\item The Torino scale\idx{Torino scale}'s two axes are not equally supported: probability is
measured and energy is inferred.
\item The relation \(\delta\ln E = 3\,\delta\ln D + \delta\ln \rho\) makes
the diameter the dominant term, so a factor of two in energy needs a
fourteen per cent diameter.
\item The Level 1 probability threshold is proportional to the inverse of the
energy, so a mis-sized object can cross a notification boundary with its
astrometry untouched.
\item The scale is a communication instrument and is not at fault; the
uncertainty enters before the scale is applied.
\end{keytakeaways}
